\let\negmedspace\undefined
\let\negthickspace\undefined
\documentclass[journal]{IEEEtran}
\usepackage[a5paper, margin=10mm, onecolumn]{geometry}
%\usepackage{lmodern} % Ensure lmodern is loaded for pdflatex
\usepackage{tfrupee} % Include tfrupee package

\setlength{\headheight}{1cm} % Set the height of the header box
\setlength{\headsep}{0mm}     % Set the distance between the header box and the top of the text

\usepackage{gvv-book}
\usepackage{gvv}
\usepackage{cite}
\usepackage{amsmath,amssymb,amsfonts,amsthm}
\usepackage{algorithmic}
\usepackage{graphicx}
\usepackage{textcomp}
\usepackage{xcolor}
\usepackage{txfonts}
\usepackage{listings}
\usepackage{enumitem}
\usepackage{mathtools}
\usepackage{gensymb}
\usepackage[breaklinks=true]{hyperref}
\usepackage{tkz-euclide} 
\usepackage{listings}
% \usepackage{gvv}                                        
\def\inputGnumericTable{}                   
\usepackage[latin1]{inputenc}                                
\usepackage{color}                                            
\usepackage{array}                                            
\usepackage{longtable}                                       
\usepackage{calc}                                             
\usepackage{multirow}                                         
\usepackage{hhline}                                           
\usepackage{ifthen}                                           
\usepackage{lscape}
\usepackage{circuitikz}
\usepackage{comment}
\tikzstyle{block} = [rectangle, draw, fill=blue!20, 
text width=4em, text centered, rounded corners, minimum height=3em]
\tikzstyle{sum} = [draw, fill=blue!10, circle, minimum size=1cm, node distance=1.5cm]
\tikzstyle{input} = [coordinate]
\tikzstyle{output} = [coordinate]
%--------------------------

\begin{document}
	
	\bibliographystyle{IEEEtran}
	\vspace{3cm}
	
	\title{10.7.90}
	\author{EE25BTECH11042 - Nipun Dasari}
	\maketitle
	
	\renewcommand{\thefigure}{\theenumi}
	\renewcommand{\thetable}{\theenumi}
	\setlength{\intextsep}{10pt}
	
	\numberwithin{equation}{enumi}
	\numberwithin{figure}{enumi}
	\renewcommand{\thetable}{\theenumi}
	
	\textbf{Question}:\\
	A circle touches the line $y=x$ at a point P such that $OP=4\sqrt{2}$, where O is the origin. The circle contains the point $(-10,2)$ in its interior and the length of its chord on the line $x+y=0$ is $6\sqrt{2}$. Determine the equation of the circle.
	
	\solution
	
	Let the point of tangency be $\vec{P}$. Since it lies on the line $y=x$, its coordinates must be equal.
	\begin{align}
		\vec{P} = \begin{myvec}{p \\ p}\end{myvec}
	\end{align}
	The distance from the origin is given as $\|\vec{P}\| = 4\sqrt{2}$.
	\begin{align}
		\norm{\vec{P}} = \sqrt{p^2 + p^2} = |p|\sqrt{2} &= 4\sqrt{2} \nonumber \\
		\implies |p| &= 4
	\end{align}
	This gives two possible points of tangency.
	\begin{align}
		\vec{P}_1 = \begin{myvec}{4 \\ 4}\end{myvec} \quad \text{and} \quad \vec{P}_2 = \begin{myvec}{-4 \\ -4}\end{myvec}
	\end{align}
	
	Let the center of the circle be $\vec{C}$.
	\begin{align}
		\vec{C} = \begin{myvec}{h \\ k}\end{myvec}
	\end{align}
	The center lies on a line perpendicular to the tangent line, which has the matrix equation $\begin{myvec}{1 & -1}\end{myvec}\vec{x} = 0$. The normal to the tangent is $\vec{n}_t = \begin{myvec}{1 \\ -1}\end{myvec}$. A normal to the line of centers, $\vec{n}_c$, must be orthogonal to $\vec{n}_t$.
	\begin{align}
		\vec{n}_c = \begin{myvec}{1 \\ 1}\end{myvec}
	\end{align}
	The equation for the line of centers is $\vec{n}_c^\top\vec{x} = d$, where $d = \vec{n}_c^\top\vec{P}$.
	For $\vec{P}_1$, the line of centers is:
	\begin{align}
		\begin{myvec}{1 & 1}\end{myvec}\vec{C} = \begin{myvec}{1 & 1}\end{myvec}\begin{myvec}{4 \\ 4}\end{myvec} = 8
	\end{align}
	For $\vec{P}_2$, the line of centers is:
	\begin{align}
		\begin{myvec}{1 & 1}\end{myvec}\vec{C} = \begin{myvec}{1 & 1}\end{myvec}\begin{myvec}{-4 \\ -4}\end{myvec} = -8
	\end{align}
	
	The circle has a chord on the line with equation $\begin{myvec}{1 & 1}\end{myvec}\vec{x} = 0$. The distance from the center $\vec{C}$ to this line is $d_{chord}$.
	\begin{align}
		d_{chord} = \frac{|\vec{n}_c^\top\vec{C}|}{\|\vec{n}_c\|} = \frac{|h+k|}{\sqrt{2}}
	\end{align}
	The radius $r$, distance $d_{chord}$, and half the chord length ($L=3\sqrt{2}$) are related by $r^2 = d_{chord}^2 + L^2$.
	\begin{align}
		r^2 = \left(\frac{|h+k|}{\sqrt{2}}\right)^2 + (3\sqrt{2})^2 = \frac{(h+k)^2}{2} + 18
	\end{align}
	The radius is also the distance from the center to the point of tangency, $r^2 = \|\vec{C}-\vec{P}\|^2$.
	
	For the first case ($\begin{myvec}{1 & 1}\end{myvec}\vec{C} = 8$ and $\vec{P}=\begin{myvec}{4 \\ 4}\end{myvec}$):
	\begin{align}
		r^2 &= \frac{8^2}{2} + 18 = 50 \\
		\norm{\vec{C}-\vec{P}}^2 &= (h-4)^2 + (k-4)^2 = (h-4)^2 + (8-h-4)^2 = 2(h-4)^2 = 50 \nonumber \\
		(h-4)^2 &= 25 \implies h-4 = \pm 5 \implies h=9 \text{ or } h=-1
	\end{align}
	This gives centers $\vec{C}_A = \begin{myvec}{9 \\ -1}\end{myvec}$ and $\vec{C}_B = \begin{myvec}{-1 \\ 9}\end{myvec}$.
	
	For the second case ($\begin{myvec}{1 & 1}\end{myvec}\vec{C} = -8$ and $\vec{P}=\begin{myvec}{-4 \\ -4}\end{myvec}$):
	\begin{align}
		r^2 &= \frac{(-8)^2}{2} + 18 = 50 \\
		\norm{\vec{C}-\vec{P}}^2 &= (h+4)^2 + (k+4)^2 = (h+4)^2 + (-8-h+4)^2 = 2(h+4)^2 = 50  \\
		(h+4)^2 &= 25 \implies h+4 = \pm 5 \implies h=1 \text{ or } h=-9
	\end{align}
	This gives centers $\vec{C}_C = \begin{myvec}{1 \\ -9}\end{myvec}$ and $\vec{C}_D = \begin{myvec}{-9 \\ 1}\end{myvec}$.
	
	The point $\vec{Q} = \begin{myvec}{-10 \\ 2}\end{myvec}$ is in the circle's interior, so $\|\vec{Q}-\vec{C}\|^2 < r^2 = 50$.
	\begin{align}
		\text{For } \vec{C}_A=\begin{myvec}{9 \\ -1}\end{myvec}: \|\begin{myvec}{-19 \\ 3}\end{myvec}\|^2 = 370 > 50 \\
		\text{For } \vec{C}_B=\begin{myvec}{-1 \\ 9}\end{myvec}: \|\begin{myvec}{-9 \\ -7}\end{myvec}\|^2 = 130 > 50 \\
		\text{For } \vec{C}_C=\begin{myvec}{1 \\ -9}\end{myvec}: \|\begin{myvec}{-11 \\ 11}\end{myvec}\|^2 = 242 > 50 \\
		\text{For } \vec{C}_D=\begin{myvec}{-9 \\ 1}\end{myvec}: \|\begin{myvec}{-1 \\ 1}\end{myvec}\|^2 = 2 < 50
	\end{align}
	The correct center is $\vec{C}_D = \begin{myvec}{-9 \\ 1}\end{myvec}$.
	
	The center is $\begin{myvec}{-9 \\ 1}\end{myvec}$ and radius squared is $r^2=50$. The equation is $(x+9)^2 + (y-1)^2 = 50$.
	\begin{align}
		x^2+y^2+18x-2y+32=0
	\end{align}
	In the matrix form $\vec{x}^\top\vec{V}\vec{x} + 2\vec{u}^\top\vec{x} + f = 0$:
	\begin{align}
		\begin{myvec}{x & y}\end{myvec} \begin{myvec}{1 & 0 \\ 0 & 1}\end{myvec} \begin{myvec}{x \\ y}\end{myvec} + 2\begin{myvec}{9 & -1}\end{myvec} \begin{myvec}{x \\ y}\end{myvec} + 32 = 0
	\end{align}
	\\
	\begin{figure}[h!]
		\centering
		\includegraphics[width=0.6\columnwidth]{figs/Figure_1.png} 
		\caption*{Fig }
		\label{Fig1}
	\end{figure}
	\begin{figure}[h!]
		\centering
		\includegraphics[width=0.6\columnwidth]{figs/Figure_2.png} 
		\caption*{Fig }
		\label{Fig2}
	\end{figure}
\end{document}